\documentclass{scrbook}
\usepackage[greek,spanish]{babel}
\usepackage[utf8]{inputenc}
\usepackage[spanish]{varioref}
\usepackage{fancyref}
\usepackage{fancybox}%recuadros
\usepackage[pdftex, breaklinks=false, colorlinks=true, linkcolor=black, anchorcolor=black, urlcolor=blue, citecolor=red, pagebackref=true]{hyperref}
\usepackage{graphicx}%graficos
\usepackage{multicol}
\usepackage{color} %color
\usepackage{xcolor} %mas opciones de color 
\usepackage{colortbl} %colorear tablas
\usepackage{lscape} % hojas apaisadas
\usepackage{cmap} % buscar en el PDF
\usepackage{listings} %para nada, al principio para listados
\usepackage{textcomp}
\usepackage{ucs} 
\usepackage{hyphenat}
\usepackage{fancyvrb} %inclusion de ficheros externo en modo verbatim

\definecolor{gris}{gray}{0.75}
\definecolor{azul}{rgb}{0,0,0.30}
\definecolor{grisclaro}{gray}{0.85}


\newenvironment{ejemplo}%
  {\VerbatimEnvironment
    \begin{Sbox}
      \begin{minipage}{\linewidth}
       \begin{Verbatim}
  }
  {\end{Verbatim}
 \end{minipage}
\end{Sbox}
\setlength{\fboxsep}{8pt}\colorbox{gris}{\TheSbox}}

\newenvironment{changemargin}[2]{
  \begin{list}{}{%
    \setlength{\topsep}{0pt}%
    \setlength{\leftmargin}{#1}%
    \setlength{\rightmargin}{#2}%
    \setlength{\listparindent}{\parindent}%
    \setlength{\itemindent}{\parindent}%
    \setlength{\parsep}{\parskip}%
  }%
  \item[]
}
{
\end{list}
}




\author{Fabián Sellés Rosa}
\title{\begin{center} \includegraphics[width=0.3\textwidth]{images/logouca} \end{center} \vfill{} \textbf{\huge{Escuela Superior de Ingeniería}} \vfill{} \textit{Ingeniería Informática} \vfill{} Sistema de deteccion de ataques usando containers \\ }
\date{}
\makeindex{}
\makeglossary{}



\begin{document}
\newcommand{\parrafo}[1]{\vspace{20pt}\underline{\textbf{#1}}}

\maketitle
%fuente LaTeX inspirado en el fuente de Antonio Garcia Dominguez en su PFC XMLEye
\thispagestyle{empty}

\includegraphics[width=0.2\textwidth]{images/logouca}

\bigskip
\bigskip
\bigskip
  
  \begin{changemargin}{3em}{3em}

    \begin{center}
      {\Huge \textsc{\nohyphens{Escuela Superior de Ingeniería}}}
      
      \bigskip
      \bigskip
      
      {\huge \nohyphens{Ingeniería Técnica en Informática de Sistemas}}
      
      \bigskip
      \bigskip
      \bigskip
      \bigskip
      
      {\LARGE \nohyphens{Análisis de implantación de telefonía IP mediante Software Libre}}
      
      \bigskip
      \bigskip
      \bigskip
      \bigskip
      
    \end{center}
  \end{changemargin}

  \begin{flushleft}
    \Large

    \textsc{Departamento}: \nohyphens{Lenguajes y Sistemas Informáticos.} \\
    \textsc{Director de proyecto}: \nohyphens{Manuel Palomo Duarte.} \\
    \textsc{Autor del Proyecto}: \nohyphens{Fabián Sellés Rosa}. \\
  \end{flushleft}
  
  \bigskip
  \bigskip
  \bigskip
  
  \begin{flushright}
    \large
    Cádiz, \today
    
    Fdo.: Fabián Sellés Rosa
    
  \end{flushright}

\bigskip
\bigskip
\bigskip
\bigskip

\newpage{}

Este documento se halla bajo la licencia FDL (Free Documentation License). Según estipula la
licencia, se muestra aquí el aviso de copyright. Se ha usado la
versión inglesa de la licencia, al ser la única reconocida
oficialmente por la \href{http://www.fsf.org/}{}FSF (Free Software Foundation).

\begin{quote}
  Copyright \copyright  2009  Fabián Sellés Rosa.
  
  Permission is granted to copy, distribute and/or modify this document
  under the terms of the GNU Free Documentation License, Version 1.2
  or any later version published by the Free Software Foundation;
  with no Invariant Sections, no Front-Cover Texts, and no Back-Cover Texts.
  A copy of the license is included in the section entitled "GNU
  Free Documentation License".
\end{quote}
\newpage{}
%%% Local Variables: 
%%% mode: latex
%%% TeX-master: "index"
%%% End: 

\thispagestyle{empty}

\begin{center}
  \begin{flushright}
    a Arturo, mi mejor amigo, también por pagar tantas cenas. \\
    \vspace{0.3cm}
    a Juanca, un buen amigo y un gran profesional aunque no se lo crea. \\
    \vspace{0.3cm}
    a Pablo, un buen amigo, profesional y alguien que aunque es más joven que yo es quién quiero ser de mayor.\\
    \vspace{0.3cm}
    a Ángel, y toda la compañía, por obligarme a descansar de vez en cuando.\\
    \vspace{0.5cm}
    a mi tutor Manuel Palomo y a Rafael Rodriguez, por su apoyo y confianza.\\
    \vspace{0.5cm}
    a los profesores y personas que me enseñaron muchas de las cosas que sé y que debía hacer, y me enseñaron como no ser.\\
    \vspace{0.5cm}
    a Stallman, por iniciar el movimiento del software libre. Y a la comunidad, por llevarlo a la práctica.\\
    \vspace{0.5cm}
    a Knuth, Floyd, Dijkstra , Turing  y tantos otros conocidos y anónimos que crearon este maravilloso mundo en el que me han dejado participar.


  \end{flushright}
\end{center}

% nada de momento

%%% Local Variables: 
%%% mode: latex
%%% TeX-master: "memoria"
%%% End: 

\tableofcontents{}
\listoffigures{}

\chapter{Introducción}

La digitalización de la sociedad no es una idea de futuro. Hace años que funciones críticas para la sociedad
como son las tareas ejercidas en hospitales, entidades financieras o centrales nucleares son gestionadas mediante sistemas informáticos. 
La ubicuidad de dispositivos, ordenadores, móviles, dispositivos electrónicos del \emph{internet of things}
(electrodomésticos, montajes con Arduino o Raspberry Pi...) unido al afianzamiento de la informatización en el ámbito empresarial, ha generado
que dichos sistemas y dispositivos se consideren activos a proteger de atacantes de índole privada o gubernamental.

La seguridad informática ha pasado de ser un área de interés para dichas corporaciones a ser considerada de interés general.
Hasta no hac tanto tiempo, hubiera sido imposible que una noticia sobre un \emph{Ransomware} apareciese en prensa o televisión no especializada y, sin embargo, esto ha ocurrido
al menos una docena de veces este 2017.

El sector de la seguridad informática es especialmente interesante desde el punto de vista del defensor, puesto que
los sistemas de detención y prevención han de cotrarrestar con talento y recursos limitados los ataques de adversarios
que possen muchos más recursos, talento y motivación para triunfar.

Es por ello que se debe asumir que -en dicho escenario- la derrota, es decir la intrusión, tiene un alto factor de éxito y es 
por ello que la gestión de riesgos no se basa en la hipótesis de si el ataque será exitoso o no, sino en que, cuando
este ataque suceda, se aprenda lo máximo posible para evitarlo en otra ocasión y limitar sus efectos.

Así pues, este proyecto se enmarca en el campo de la contrainteligencia y en él se describe un método que permite al defensor conocer cómo algunos atacantes intentan vulnerar sus sistemas. 






%%% Local Variables: 
%%% mode: latex
%%% TeX-master: "index"
%%% End: 

\chapter{Descripción General}

\section{Telefonía y redes IP}

\subsection{Generalidades}
\label{sec:generalidades}


En el presente capitulo se presenta una introducción de conceptos fundamentales para entender el funcionamiento de la telefonía IP.


La telefonía existe desde hace ya más de 100 años, las necesidades de comunicación y las nociones de distancia han cambiado con ella y a causa de ésta. Desde la aparición y sobre todo la popularización de Internet, otros medios más ligados a la informática como el e-mail o la mensajería instantánea parecen haber sustituido a la telefonía en determinadas circunstancias.


Sin embargo, hay organizaciones o circunstancias que nos hacen depender del teléfono. A pesar de la irrupción de microportatiles y PDAs, no es común encontrarse a mano un ordenador con conexión a Internet y hay comunicaciones que necesitan más interactividad que la que proporciona cruzar un par de mensajes de texto.

Así nos encontramos que es habitual mantener los dos sistemas no relacionados sin conectarlos aunque, sin embargo, comparten la misma finalidad: la necesidad de transporte e intercambio de información.


\subsection{Caracterización de la voz humana}

La voz humana es una señal analógica no periódica cuyo rango de operación en frecuencias está entre los 20 y los 20000 Hz. Sin embargo, su ancho de banda útil, donde se concentra la mayor información, se encuentra entre los 300 y los 4000 Hz.

La rápida implantación de terminales hizo necesaria la creación de centrales (primero operadas por humanos y luego automáticas) y la creación de protocolos de telefonía.

Sin entrar en detalles de cada uno de los protocolos hay ciertas señales (el número que marcamos, el colgado, el descolgado, u otras más complejas) que deben ser enviadas a la central para su tratamiento. Para esto existen dos alternativas:

\begin{list}{-}{}
  \label{lst:tipos-de-sel}
\item Señalización \emph{in-band} (dentro de banda), donde las señales viajan junto a la voz. Esto permite que el usuario pueda llegar a conocerlas e incluso a emularlas (los inicios del phreaking están muy relacionados con esto,una caja de cereales y un silbato capaz de emular los tonos de la central, veasé \cite{wikipedia-crunch}).

  \begin{figure}[!b]
    \centering
    \includegraphics[scale=0.5]{images/silbato_crunch}
    \caption{Silbato capaz de emular un tono para manipular las antiguas centrales telefónicas.}
  \end{figure}


 También se puede enviar la señalización junto a la voz de manera oculta al usuario, utilizando cabeceras que se generan en el terminal, que posteriormente son captadas por la central y eliminadas cuando se reenvía para el terminal de destino (encapsulación). En telefonía digital, puede hacerse un proceso similar generando bits de guarda que marcan el inicio de una señal que son recibidos en la central telefónica, procesados y eliminados antes de su reenvío al destinatario. (Veasé \cite{wikipedia-inband})

\item Señalización \emph{out-band} (fuera de banda), se define un canal separado donde viaja la señalización. En las lineas RDSI por ejemplo, en un acceso básico existen dos canales de datos cada uno de 64 Kib  y uno de señalización de 16 Kib (2B+D), la voz viaja por uno de los canales B y toda la señalización por el canal D.

 \end{list}

La señalización es relevante para la central o agente intermedio, ya que determina qué está ocurriendo realmente entre usuarios. Veasé la figura \ref{fig:modelo-funcional}

\begin{figure}
  \centering
  \includegraphics[width=0.5\textwidth]{images/modelo_telefonia}
  \caption{Modelo funcional de Telefonía}
  \label{fig:modelo-funcional}
\end{figure}

\subsection{Codecs}
\label{sec:codecs-1}

Para poder enviar la voz mediante medios digitales, debemos \emph{discretizarla}. Para este proceso de \emph{discretización o cuantificación} se utiliza un codec\footnote{\emph{cod}ificador-\emph{dec}odificador}. Un codec es una \emph{caja negra} física o lógica, que en la entrada recibe una señal de audio y a la salida devuelve una cadena de bits, o viceversa.

Los parámetros habituales para comprobar las características de un codec son:

\begin{description}
\item[Frecuencia de muestreo], según el Teorema de Nyquist\footnote{o Nyquist-Shannon} para poder reproducir fielmente una señal con frecuencia máxima $f$ es necesario tomar muestras con una frecuencia $2f$. En telefonía tradicional y en VoIP la frecuencia de muestreo suele estar en los 8000 Hz, pues se considera que la mayor parte de los tonos de voz se concentran hasta los 4000 Hz. No obstante hay codecs que utilizan frecuencias más altas siendo 44100 Hz la máxima (pues el oído humano no escucha frecuencias superiores de 22,5 KHz).
\item[Número de bits por muestra]. al tomar muestras para discretizar una señal continua, utilizar más bits hace que la señal reconstruida ajuste mejor a la original.
\item[Perdida] la codificación desprecia armónicos no audibles o no importantes en la señal de audio original, generando una señal que una vez reproducida no es \emph{completamente} fiel a la de la entrada pero si \emph{suficientemente} fiel.
\item[Número de canales] un sólo canal (mono), dos (estéreo) o más. En telefonía suelen utilizarse codecs que utilicen un único canal, debido a que más canales requeriría más ancho de banda para su transmisión.
\end{description}



\subsection{Redes IP}
\label{sec:redes-ip}

IP (\emph{Internet Protocol}) es un protocolo de entrega de paquetes de máximo esfuerzo no confiable. Las redes conmutadas de paquetes, como las que utilizan el protocolo IP, difieren de otras redes como las conmutadas por circuitos o por tramas. La telefonía tradicional ha venido utilizando conmutación por circuitos, cuando se utiliza este tipo de red de comunicación entre llamador y llamante se establece un circuito físico o lógico que se puede seguir directamente desde el origen al destino de la llamada.

En el caso de las redes conmutadas de paquetes la información que se desea transmitir se trocea en paquetes, los paquetes se envían desde origen a destino. Pero hay que hacer algunas matizaciones: dos paquetes de una misma \"conversación\" no tienen porque seguir la misma ruta y el orden en que lleguen los paquetes al destino depende del estado de la red, y la latencia de la ruta escogida por cada uno.

IP es un protocolo que está dentro de la pila TCP/IP, que es un conjunto de protocolos que cubren \footnote{Entre ellos HTTP , DNS , ARP , ICMP \ldots} todos los servicios de red en una red de paquetes. Si el lector no tiene una solida formación en redes o en TCP/IP, una buena referencia es \cite{Stallings07}.

Los protocolos de señalización y transporte a tiempo real del siguiente capitulo utilizan la pila TCP/IP, eso quiere decir que todos los protocolos acabarán utilizando un protocolo de transporte como TCP o UDP en combinación con un protocolo de entrega de paquetes como IP.



\subsection{Protocolos abiertos para la señalización}

Un protocolo abierto es un protocolo del cual podemos leer su especificación. Esto es importante porque nos permite crear nuestros propios clientes y servidores, lo que nos libera de una implementación concreta de un fabricante u proveedor.

Una de las primeras cosas que hay que tener en cuenta al tratar un protocolo es que la especificación debe de materializarse en una implementación, y comprobar que una implementación dada responda adecuadamente a lo estipulado en la especificación.

Sin embargo, \emph{podemos} (si lo necesitamos) comprobar como de fiel es a la especificación una implementación dada por los mensajes que esta devuelve. Esto no es posible con protocolos cerrados o propietarios donde debemos confiar en la palabra del fabricante u proveedor.




\section{Protocolos abiertos de VoIP}
\label{sec:prot-abiert-de}


\subsection{H323}

H323 fue el primer protocolo abierto de telefonía IP. Surge como una evolución \" natural \" de otros protocolos de videoconferencia y telecomunicaciones por parte de la ITU-T\footnote{The Telecommunication Standardization Sector (ITU-T)}.Es parte de la serie de protocolos H.32x que es un conjunto de estándares que definen los elementos necesarios para la Telefonía a través de IP.

En la familia H.32X hay una serie de figuras definidas, cada figura cumple con una serie de atribuciones:  
\begin{description}

\item[Terminal] dispositivo para iniciar o recibir la comunicación.
\item[Gatekeeper] (\emph{controlador de dominio}) dispositivo que proporciona:
  \begin{enumerate}
  \item Traducción de direcciones
  \item Control de admisiones
  \item Control de Ancho de banda
  \item Administración de Zonas
  \item Control de señalización de llamada
  \item Autorización de llamadas
  \item Administración de ancho de banda
  \item Administración de llamadas
  \end{enumerate}

\item[Gateway] (\emph{convertidor de medios}) dispositivo que permite la conversión de redes H323 a otras redes como PSTN (RTC), ISDN (RDSI), etc 
\end{description}

A continuación, se muestra un listado de los principales protocolos de la familia H.32X:

\begin{list}{-}{}
\label{lst:h32x}
\item H.225.0 Registro, Admisión y eStatus (RAS), utilizado entre un terminal H.323 y un Gatekeeper para proporcionar resolución de direcciones y servicios de control de admisión. 

\item H.225.0 Señalización de llamada, utilizado entre dos entidades H323 para establecer la comunicación (terminal $\leftrightarrow $ terminal , Gatekeeper $\leftrightarrow $ terminal , \ldots ).
\item H.245 protocolo de control para comunicaciones multimedia. Estas describen los mensajes y procedimientos usados para el intercambio de capacidades, apertura y cierre de canales lógicos para audio, vídeo y datos, controles e indicaciones.

\item Real-time Transport Protocol (RTP), el cual es usado para enviar y recibir información multimedia (voz , vídeo , o texto) entre dos entidades.
\item H.235 describe la seguridad en H.323, para señalización y multimedia.

\item H.239 describe el canal dual usado en videoconferencia, usualmente una para vídeo en tiempo real y otro para imágenes fijas.

\item H.450 describe varios servicios suplementarios.

\item H.460 define extensiones opcional que pueden ser implementados en un terminal o Gatekeeper, especialmente para soslayar las NAT/Firewall.

\end{list}

El anterior listado nos da una idea de la complejidad del protocolo H323, pues involucra a varios protocolos que cooperan para ofrecer un servicio. 


El acceso a la especificación de los protocolos H32X no es gratuito, el estándar se ha de comprar, siguiendo la política de distribución de estandarización ISO. Aún así, el ITU-T ha publicado algunos estándares de manera libre que pueden verse en \cite{H323}.

El proceso de llamada es bastante complejo en H323, y la petición de servicios no es ortogonal, es cambiante según las figuras que estén presentes. El siguiente es un ejemplo de llamada entre dos terminales H323 (Figura \ref{fig:llamada-H323}).


\begin{figure}
  \centering
  \includegraphics[width=0.5\textwidth]{images/H323}
  \caption{Llamada entre dos terminales H323}
  \label{fig:llamada-H323}
  \footnote{Imagen reproducida con permiso del autor , \emph{Fuente : http://www.voipforo.com/images/H323-comunicacion.gif}}
\end{figure}

\subsection{SIP}


SIP es un protocolo que define cómo establecer sesiones multimedia entre dos usuarios. En él no se explica ni se define como pasar la voz, sólo se define como señalizarlo. SIP es un protocolo completo, por lo que al hablar de él en realidad nos estamos refiriendo a una familia de protocolos que SIP utiliza o que extienden a éste.

En concreto en el propio RFC\cite{rfc3261}, se define: 
\begin{quotation}
  

\emph{"SIP es un protocolo de control de la capa de aplicación que puede establecer, modificar, y terminar sesiones multimedia como llamadas telefónicas por Internet. SIP puede también invitar participantes a sesiones existentes, como conferencias."
}
\end{quotation}

SIP soporta cinco funciones de las comunicaciones multimedia:
\begin{list}{$\bullet$}{}
\item Localización del usuario: determina el lugar del usuario final para la comunicación.
\item Disponibilidad del usuario: determina la voluntad del llamado de unirse a una comunicación.
\item Capacidades del usuario: determina los medios que se usarán.
\item Sesión de configuración: "llamar" , establecimiento de los parámetros de sesión para llamante y llamado.
\item Gestión de la sesión: incluyendo transferencia y terminación de las sesiones, modificación e invocación de servicios.
\end{list}


A diferencia de H323, SIP  esta disponible de manera publica y gratuita en la web de los RFC. (\cite{rfc3261})

SIP esta basado en el protocolo HTTP \cite{rfc2616}, del cual toma el modelo general y algunos códigos de respuesta. Cualquier paquete SIP se escribe en texto plano, y puede observarse mediante cualquier analizador de red (Veasé sección \vref{sec:herr-de-anal}).

\subsubsection{Ejemplo de conversación con SIP}
\label{sec:ejemplo}

El ejemplo más básico es el que involucra dos proxys y dos terminales, que forman el denominado trapezoide SIP (Figura \ref{fig:trapezo-sip}). Ésta es la topología de red más usual y sencilla, en la que intervienen dos terminales conectados cada uno de ellos a un servidor proxy SIP.

\begin{figure}[b]
  \centering
  \includegraphics[width=0.6\textwidth]{images/trapezoide}
  \caption{Trapezoide SIP}
\label{fig:trapezo-sip}
\end{figure}

Si \emph{tlf1} quisiese llamar a \emph{tlf2}, el diagrama \ref{fig:conv-sip} nos muestra la secuencia de mensajes que se daría. Para cada mensaje se indica el orden en que se emite, el tipo de mensaje y la direccionalidad del mismo. Hay que tener en cuenta que en principio \emph{tlf1} no puede llegar directamente ni a \emph{tlf2} ni a \emph{proxy2.com}, así que deja esa responsabilidad a su servidor SIP, que si lo conoce (en caso contrario le respondería con un mensaje 404 de no encontrado\footnote{¿Le suena?}).

\begin{figure}[!h]
  \centering
 \label{fig:conv-sip} 
 \includegraphics[width=0.8\textwidth]{images/conversacionSIP}
  \caption{Conversación SIP usual de \emph{tlf1} a \emph{tlf2}}

\end{figure}

Dentro de este pequeño ejemplo podemos concretar los elementos de cualquier conversación SIP, los pasos del 1 al 12 y del 13 al 14, son mensajes de señalización que no involucran audio, los mensajes de 1 a 3 realizan la invitación, los mensajes del 6 al 8 muestran al llamante que se ha realizado la llamada, y los mensajes del 9 al 11 son la aceptación de la llamada por el llamado. Una vez realizado todo esto, es cuando realmente se procede a la comunicación entre terminales.

Nótese que la comunicación entre terminales se produce en el mensaje 12, hasta entonces la comunicación es a través de servidores SIP. Pero una vez que se acepta la llamada la comunicación es ya directa entre terminales, no se notifican a los servidores los mensajes.

\subsubsection{Transacciones y diálogos, el modelo de señalización SIP}
\label{sec:trans-y-dial}

En el ejemplo anterior, veíamos algunos actores de una comunicación entre dos terminales SIP y dos servidores proxy. SIP es un modelo de petición/respuesta, a una petición pueden seguirle una o más respuestas. Definimos así:

\begin{quotation}
\emph{Una \textbf{transacción} se compone de una petición inicial seguida de una o varios mensajes de respuesta. E.j.: Un mensaje INVITE comienza una transacción que acaba tras varias respuestas con un código 200.}
\end{quotation}

En cambio, un diálogo es la comunicación que se produce entre dos mismos agentes de usuario, i.e. entre dos clientes o entre un cliente y un servidor. 

\begin{quotation}
\emph{Un \textbf{diálogo} es una o más transacciones que se producen entre dos mismos agentes de usuario.Para identificar el diálogo se miran las cabeceras} \verb|From:| \emph{,} \verb|To:| \emph{y } \verb|Call-ID:|. \emph{el orden de mensajes de un diálogo se determina con la cabecera} \verb|Cseq:|
\end{quotation}

\label{sec:agente-de-usuario}

Por su parte un \textbf{agente de usuario}, puede dividirse en dos partes, un \textbf{servidor de agente de usuario} que recibe las peticiones de otros agentes de usuario clientes, y un \textbf{agente de usuario cliente} que realiza las peticiones a otros agentes de usuario servidores. Un terminal se compone de un agente de usuario servidor y un agente de usuario cliente, un servidor SIP se compone de uno o varios agentes de usuario servidor y cliente.

\subsubsection{Peticiones y respuestas}
\label{sec:peticiones-sip}

Una petición SIP indica el usuario o servicio al que se quiere acceder. Una petición tiene la siguiente estructura:

\begin{verbatim}
Metodo dirección-adonde-la-peticion Version-SIP CRLF
\end{verbatim}

Donde método, puede ser uno de los seis siguientes, definidos en el estándar, o algún otro, definido en las extensiones para SIP\footnote{como PUBLISH para notificar un cambio de evento\cite{rfc3903}, o MESSAGE\cite{rfc3428} para mensajeria instantanea.}:

\begin{description}
\item[INVITE], para comenzar una sesión con otro usuario.

\item[REGISTER], para registrar la localización actual del usuario y poder comenzar a utilizar los servicios SIP.
\item[ACK], para confirmar los mensajes que se produzcan en una sesión.
\item[CANCEL] , para denegar el establecimiento de una sesión. (para colgar, por ejemplo)
\item[BYE], para acabar una sesión establecida.
\item[OPTIONS], para preguntar al servidor sobre las capacidades del mismo.

\end{description}

Cada una de estas peticiones genera respuestas, las respuestas generadas siguen el siguiente formato:

\begin{verbatim}
VERSION SIP Codigo-estado razon-de-respuesta CRLF
\end{verbatim}

Donde código de estado es un código numerico que indica el tipo de respuesta, y la razón de respuesta es una razón textual que explica más detalladamente la respuesta. El primer dígito del código numerico marca el significado de la respuesta, pudiendo ser:

\begin{list}{$\bullet$}{}
\item 1xx: Provisional, petición recibida, continua el procesamiento de la petición. E.j. el 100 significa \"intentando\".
\item 2xx: Éxito, la petición fue recibida,entendida y aceptada.
\item 3xx: Redirección, se necesitan más acciones para completar la petición, o el servidor proxy no puede acabar la petición y se reencamina a otro que si pueda.

\item 4xx: Error del cliente, la petición esta mal formada o no puede ser procesada en ese servidor.

\item 5xx: Error del servidor, el servidor fallo al procesar una petición aparentemente valida.
\item 6xx: Fallo global, la petición no puede ser procesada en ningún servidor.
\end{list}

\subsubsection{Tipos de servidores SIP}
\label{sec:tipos-de-servidores}

El RFC (\cite{rfc3261}) define una serie de entidades lógicas que tienen funciones muy concretas: son las encargadas de cumplir las funciones del protocolo.

\begin{itemize}


\item \parrafo{Servidor registrar}

\begin{figure}[!h]
\label{fig:sip-registrar}  
\centering
  \includegraphics[width=0.6\textwidth]{images/registrar}
  \caption{Servidor registrar}
\end{figure}

\emph{Un servidor registrar es el} servidor que recibe los mensajes \verb|REGISTER| , es el encargado de pedir la autenticación al usuario y de almacenar su localización actual (ip,nombre de usuario y puerto) en una \emph{base de datos de localizaciones}. Cuando otros servidores (como un servidor proxy) busquen , por ejemplo, al usuario sip:tlf@proxy1.com , la base de datos de localizaciones le devolverá \\ \emph{sip:tlf1@a.b.c.d:puerto}. Es un servidor lógico que generalmente se utiliza un servidor proxy. 


\item \parrafo{Servidor redirect}

\begin{figure}[!h]
  \centering
  \includegraphics[width=0.5\textwidth]{images/redirect}
  \caption{Servidor redirect}
  \label{fig:sip-redirect}
\end{figure}

\emph{Un servidor redirect} indica al cliente la localización actual de la persona a la que llama, pero no inicia la comunicación. En el ejemplo del trapezoide \ref{sec:ejemplo} , tras el mensaje 1, le seguiría una respuesta del servidor redirect con un mensaje de respuesta 3XX, donde se le indica la localización del usuario que se buscaba. Entonces, el llamante debe generar un nuevo INVITE al usuario encontrado. 



\item \parrafo{Servidor proxy o SIP router}

\emph{Un servidor proxy es} el encargado de \textbf{\emph{encaminar}} una petición de sesión de un usuario al proxy más cercano que pueda dar curso a esa petición. Es también el encargado de la autenticación y de el \emph{accounting\footnote{No existe una buena traducción sin ambiguedades, el accounting es la gestión de usuarios en el sentido de permisos, servicios que puede acceder y para facturación.}.}Un ejemplo de servidor proxy funcionando se ha mostrado ya mediante el trapezoide SIP (figura \ref{fig:trapezo-sip}).

\end{itemize}
\subsubsection{Nivel de implantación de SIP}
\label{sec:implantacion-de-sip}

SIP es un protocolo estándar aprobado por el IETF, que tiene muchos dispositivos compatibles en el mercado. Es con diferencia el protocolo que más se implementa en los dispositivos de telefonía IP.


\subsubsection{XMPP}
\label{sec:xmpp}

XMPP (\emph{e\textbf{X}tensible \textbf{M}essaging and \textbf{P}resence \textbf{P}rotocol},Protocolo extensible de mensajería y comunicación de presencia) es un protocolo abierto , estandarizado por el IETF (\cite{rfc3920} \cite{rfc3921} y \cite{rfc3922}) pero con muchas extensiones en proceso de estandarización. XMPP es un protocolo que fue diseñado para mensajería instantánea, para este propósito intercambia entre clientes ficheros XML. Actualmente se está desarrollando una extensión para XMPP(Jingle \cite{jingle-overview}), que permita el flujo de audio mediante RTP de modo que Jingle se encargue de la señalización. Además el protocolo se ha diseñado intentando mantener cierta compatibilidad con SIP.

Pero además de su posible inclusión como nuevo protocolo de VoIP, XMPP tiene su propio peso en el mundo de la VoIP como protocolo de presencia. Aunque SIP tiene SIMPLE, hay pocas implementaciones del protocolo, en cambio XMPP lleva siendo implantado de manera masiva (GTalk, el cliente de mensajería instantánea de Google utiliza XMPP). 

\subsection{Transporte en tiempo real y redes IP}
\label{sec:transporte-en-tiempo}

Al transmitir datos que no necesitan llegar en tiempo real, la latencia no nos afecta demasiado, siempre y cuando los paquetes lleguen al destino. Al transmitir audio,vídeo o cualquier sesión multimedia esta situación no puede producirse, pues en este tipo de datos la temporización es muy importante para poder comprender el contenido del mismo. Así que no nos basta con cerciorarnos de que los paquetes lleguen sino que además necesitamos que lleguen en el tiempo necesario y se ordenen de la manera adecuada.

El protocolo encargado de la temporización y del reordenado de este tipo de paquetes es RTP, éste sin embargo, confía en que los paquetes llegarán en un tiempo \emph{adecuado} para que el destino este pueda reordenarlos. 

Los paquetes UDP son más pequeños que los TCP, y aunque no sean confiables, se suele utilizar este transporte para los servicios de telefonía IP. Esto unido a la creación de un \emph{jitter-buffer} (Veasé \ref{def:jitter-buffer})  y a la utilización de equipos con QoS \footnote{Calidad de servicio. Capacidad de los equipos de red de priorizar los paquetes según su relevancia}, permiten mantener controlada la latencia de red y hacer que la transmisión de datos multimedia sea lo suficientemente \emph{fluida} para humanos, sin los cortes producidos por paquetes que llegan fuera de tiempo. 

\subsection{RTP y RCTP}
\label{sec:rtp}

RTP (\emph{Realtime Transport Protocol }) y RTCP (\emph{Realtime Transport Control Protocol}), son protocolos destinados al transporte de flujos multimedia. Mientras RTP se encarga del transporte propiamente dicho, RTCP utiliza RTP para ofrecer calidad de servicio y otras funciones relacionadas con conferencias y sesiones multimedia.

\subsection{RTP y  NAT}
\label{sec:rtp-y-las}

Como ya hemos comentado, H323 y SIP son protocolos de señalización, que utilizan para el transporte multimedia RTP. Esto introduce un problema. Por ejemplo, en SIP, la señalización utiliza un único puerto, generalmente el 5060, para la administración de la red solo es necesario dejar paso a ese puerto.

Sin embargo una conversación puede tener varios flujos RTP, varios canales de audio y vídeo, y un único cliente puede utilizar varios flujos RTP (con otra cuenta de VoIP, o para streaming). RTP funciona en cualquier puerto UDP disponible por encima de 5000\footnote{Los puertos por debajo de 1024 inclusive, son los denominados puertos privilegiados. En algunos sistemas Unix además los puertos entre el 1024 y el 5000 son asignados por el sistema operativo y no por las aplicaciones de usuario.} (\cite{rfc3551}, \cite{rfc3550}). El puerto seleccionado ha de ser par, puesto que el puerto siguiente impar sera el número de puerto utilizado por el protocolo RTCP. \cite{rfc4961} . 

\subsubsection{El problema RTP-NAT}
\label{sec:el-problema}

El problema proviene cuando existe entre los dos terminales, un dispositivo que está dentro de la ruta de transporte de la red que hace NAT  o es un firewall.

Más arriba se ha comentado que RTP utiliza puertos aleatorios por encima del 5000. Al no tener un rango de puertos definido es difícil darles cabida en las reglas de filtrado de firewall o en las reglas de enrutamiento NAT para los dispositivos NAT:

\begin{figure}[!h]
  \centering
  \includegraphics[width=0.7\textwidth]{images/NAT-problema}
  \caption{El problema}
  \label{fig:NAT-Problema}
\end{figure}
Por si esto fuera poco, la implementación de la función de NAT depende del fabricante y del dispositivo. Se pueden identificar varios tipos de NAT.

\begin{description}
\item[Full Cone NAT] 

Una NAT de tipo \emph{Full Cone} es una donde todas las peticiones de las misma IP interna y puerto son redirigidas a una dirección IP y puerto. Además cualquier equipo externo puede enviar un paquete a la dirección interna utilizando la IP redirigida y el puerto. Si no se escoge ninguna redirección cada paquete proveniente de una dirección interna utiliza una redirección común con cualquier puerto. Veasé figura 

\begin{figure}[!h]
  \centering
  \includegraphics[width=0.7\textwidth]{images/Full_Cone_NAT}
  \caption{\emph{NAT Full Cone} , Fuente : \emph{Wikipedia}}
\label{fig:fnat}
\end{figure}

\item[Restricted Cone NAT] 

Todas las peticiones de una IP interna son redirigidas a una IP pública y puerto. Cualquier equipo externo sólo podrá enviar un paquete a la IP interna si ésta le envío anteriormente un paquete al equipo externo sin importar el numero de puerto. 

\begin{figure}[!h]
  \centering
  \includegraphics[width=0.7\textwidth]{images/Restricted_Cone_NAT}
  \caption{\emph{Restricted Cone} , Fuente : \emph{Wikipedia}}
  \label{fig:rnat}
\end{figure}


\item[Port Restricted Cone NAT] 

La \emph{Port Restricted Cone NAT} es similar a la \emph{Restricted Cone NAT} excepto que un equipo externo sólo podrá enviar un paquete a un equipo interno, si el equipo interno envío un paquete anteriormente a la dirección IP $X$ del equipo externo en el puerto $P$ y esté responde a la dirección interna utilizando \textbf{el mismo puerto}. Es decir, sólo se aceptara un paquete externo dirigido hacia una dirección interna y puerto si ésta envió a la dirección externa un paquete.


\begin{figure}[!h]
  \centering
  \includegraphics[width=0.7\textwidth]{images/Port_Restricted_Cone_NAT}
  \caption{\emph{Port Restricted Cone} , Fuente : \emph{Wikipedia}}
  \label{fig:rpnat}
\end{figure}

\item[Symmetric NAT] 
\end{description}

En la \emph{Symmetric NAT} una IP interna es redirigida a una dirección IP externa y un puerto. Si el mismo equipo interno envía un paquete desde la misma dirección y puerto pero a otro destino, se utiliza una redirección diferente. Es decir, para elegir una nueva redirección se utiliza un trío \emph{(dirección origen, puerto, dirección destino)} si alguno de los valores del trío cambia se utiliza una nueva redirección.

\begin{figure}[!h]
  \centering
  \includegraphics[width=0.7\textwidth]{images/Symmetric_NAT}
  \caption{\emph{Symmetric NAT} , Fuente : \emph{Wikipedia}}
  \label{fig:snat}
\end{figure}

\subsection{Soluciones al problema}
\label{sec:soluc-al-probl}

Para afrontar este problema hay algunas soluciones. Todas estas soluciones han sido concebidas para SIP o RTP, aunque pueden utilizarse para otras aplicaciones que adolezcan de los mismos problemas.

\subsubsection{Evitarlas}
\label{sec:evitarlas}

Lo más directo, aunque no es en realidad una solución, consiste en confinar las comunicaciones VoIP dentro de la LAN, y para aquellos casos en que se necesite sortear un dispositivo con NAT utilizar un túnel VPN\footnote{\emph{Virtual Private Network}, red privada virtual}.

\begin{figure}[!h]
  \centering
  \includegraphics[width=0.8\textwidth]{images/NAT-vpn}
  \caption{Utilizar túneles VPN}
  \label{fig:NAT-vpn}
\end{figure}

\subsubsection{STUN}
\label{sec:servidor-stun}

STUN\footnote{\emph{Simple Traversal of UDP over NAT}} es un protocolo que sirve para que el cliente conozca qué tipo de NAT está utilizando y su IP pública de modo que pueda intentar modificar las cabeceras para paliar este comportamiento.

\begin{figure}[!h]
  \centering
  \includegraphics[width=0.8\textwidth]{images/NAT-STUN}
  \caption{Funcionamiento de STUN}
  \label{fig:STUN-funcionamiento}
\end{figure}

Para lograr determinar el tipo de NAT, STUN utiliza el siguiente algoritmo expuesto en la figura \ref{fig:STUN-Algoritmo} :

\begin{figure}[!h]
  \centering
  \includegraphics[width=0.8\textwidth]{images/ALGORITMO-STUN}
  \caption{Algoritmo del protocolo STUN, \emph{Fuente: Wikipedia}}
  \label{fig:STUN-Algoritmo}
\end{figure}

STUN esta limitado para protocolos que utilicen UDP y tengan una determinada configuración NAT. Por lo tanto es una parte de la resolución del problema pero no lo resuelve en todos los casos.

\subsubsection{TURN}
\label{sec:servidor-turn}

TURN\footnote{\emph{Traversal using Relay NAT}} es una extensión de STUN, se encarga de retransmitir los paquetes desde el cliente al servidor y viceversa.

\begin{figure}[!h]
  \centering
  \includegraphics[width=1\textwidth]{images/NAT-TURN}
  \caption{Funcionamiento de TURN}
  \label{fig:TURN-funcionamiento}
\end{figure}

Además de los problemas de seguridad inherentes a esta solución (deben utilizarse mecanismos de autentificación y protección) el reenvío introduce una redundancia y un coste para la red y el ancho de banda importante, al tratarse de datos multimedia.


\section{PBX para telefonía IP}
\label{sec:pbx}

\subsection{Perspectiva}
\label{sec:perspectiva}

Las PBX (\ref{def:PBX}) han tenido una importancia crucial en los sistemas de telefonía tradicionales en empresas y otras organizaciones. La funcionalidad mínima que se espera de ellas es la de comunicar dos telefonos internos y conectar estos telefonos internos con una línea exterior. 

Sin embargo, existen otras funciones que con el tiempo se han implementado como un \emph{núcleo} de funciones básicas para una PBX. Generalmente una PBX utiliza un protocolo de comunicación con sus extensiones y con sus lineas externas, es posible que una misma PBX tenga varias interfaces distintas para poder ser multiprotocolo.

Una PBX es más una entidad lógica que un dispositivo fisico, pueden encontrarse PBX en hardware especifico, o PBX desarrolladas en software que se instalan en un PC generico, a estas ultimás se les denomina \emph{soft PBX}.

\subsection{¿Soft PBX?}
\label{sec:centralita-basada-en}

Historicamente las PBX han pasado de ser operadoras humanas que hacian la interconexión del cableado mediante latiguillos a circuitería analógica mediante relés, posteriormente pasando a circuitos digitales de aplicación a especifica, hasta adoptar microcontroladores y microprocesadores con sistemás operativos y software especifico a sistemás operativos "generalistas".


%%% Local Variables: 
%%% mode: latex
%%% TeX-master: "memoria"
%%% End: 





\chapter{Conclusiones y trabajo futuro}
\label{cha:concl-y-trab}

\section*{Sobre los objetivos propuestos}

Tras el trabajo realizado, podemos evaluar si se han cumplido los objetivos propuestos inicialmente. Creemos que estos han sido satisfechos, debido a que el capitulo 2 de este documento, es un punto de partida para cualquier instalador de telefonía IP para comprender los aspectos al menos más básicos de la tecnología. No existía documentación que relacionase anteriormente todos los aspectos que el neófito debe abordar (protocolos de señalización, transporte, aplicaciones disponibles, \ldots)

Así mismo se ha generado una \textbf{ documentación libre y de calidad} que no existia. Explicando, de manera detallada, la instalación y configuración de una centralita \emph{Asterisk}. Se han proporcionado ejemplos de uso (pasando desde una pequeña PYME a una organización grande como es la propia Universidad) que aunque básicos, cubren las necesidades que mediante otras tecnologías tienen implantadas muchas organizaciones.

\section*{Sobre telefonía IP y las aplicaciones de software libre}


\begin{itemize}
\item Las aplicaciones de software libre son lo suficientemente maduras como para implantar un sistema de telefonía a cualquier nivel. Desde pequeñas organizaciones a proveedores telefónicos. La flexibilidad, potencia de estas aplicaciones marcan un cambio radical con respecto al desarrollo de sistemas telefónicos clásicos.
\item SIP ha sido y va a ser por el momento el protocolo de señalización imperante. No obstante, la evolución de XMPP, y las otras aplicaciones del protocolo para mensajería instantánea y presencia, hacen que sea un protocolo a tener en cuenta a medio plazo.
\item Es necesario la creación de interfaces gráficos o mejorar los existentes para que el operador pueda gestionar su instalación. Actualmente las aplicaciones de gestión o estos interfaces están orientados al instalador, permitiendo el control de todo el sistema.
\item Debido a la aparición de terminales móviles con mejores recursos, y a la aparición de terminales con un sistema operativo propio libre (como \emph{Android}),es muy necesaria la creación de un softphone libre para móviles.
\end{itemize}


\section*{Trabajo futuro}

\begin{itemize}
\item Mejorar o crear interfaces gráficas de usuario y/o aplicaciones para \emph{Asterisk} y \emph{openSER} orientadas al operador que explote el sistema. Por ejemplo, ayudando a la creación de IVR de manera gráfica mediante \emph{Asterisk} o integrando estas herramientas en herramientas de monitorización como \emph{Nagios}.
\item Desarrollar nuevos softphones para móviles y softphones libres con interfaces gráficas más pulidas que las existentes, utilizando para ello librerias libres como \emph{PJSIP} \cite{devel-pjsip} y Sofia \cite{devel-sofia}.
\item Explorar las posibilidades como centro de domótica de \emph{Asterisk} utilizando pequeñas placas como \emph{Arduino}.
\end{itemize}

%%% Local Variables: 
%%% mode: latex
%%% TeX-master: "memoria"
%%% End: 

\appendix
%---------------------------------------------------------------------


 \begin{center}

       Version 1.2, November 2002


 Copyright \copyright 2000,2001,2002  Free Software Foundation, Inc.
 
 \bigskip
 
     51 Franklin St, Fifth Floor, Boston, MA  02110-1301  USA
  
 \bigskip
 
 Everyone is permitted to copy and distribute verbatim copies
 of this license document, but changing it is not allowed.
\end{center}


\begin{center}
{\bf\large Preamble}
\end{center}

The purpose of this License is to make a manual, textbook, or other
functional and useful document ``free'' in the sense of freedom: to
assure everyone the effective freedom to copy and redistribute it,
with or without modifying it, either commercially or noncommercially.
Secondarily, this License preserves for the author and publisher a way
to get credit for their work, while not being considered responsible
for modifications made by others.

This License is a kind of ``copyleft'', which means that derivative
works of the document must themselves be free in the same sense.  It
complements the GNU General Public License, which is a copyleft
license designed for free software.

We have designed this License in order to use it for manuals for free
software, because free software needs free documentation: a free
program should come with manuals providing the same freedoms that the
software does.  But this License is not limited to software manuals;
it can be used for any textual work, regardless of subject matter or
whether it is published as a printed book.  We recommend this License
principally for works whose purpose is instruction or reference.


\begin{center}
{\Large\bf 1. APPLICABILITY AND DEFINITIONS}
\end{center}

This License applies to any manual or other work, in any medium, that
contains a notice placed by the copyright holder saying it can be
distributed under the terms of this License.  Such a notice grants a
world-wide, royalty-free license, unlimited in duration, to use that
work under the conditions stated herein.  The \textbf{``Document''}, below,
refers to any such manual or work.  Any member of the public is a
licensee, and is addressed as \textbf{``you''}.  You accept the license if you
copy, modify or distribute the work in a way requiring permission
under copyright law.

A \textbf{``Modified Version''} of the Document means any work containing the
Document or a portion of it, either copied verbatim, or with
modifications and/or translated into another language.

A \textbf{``Secondary Section''} is a named appendix or a front-matter section of
the Document that deals exclusively with the relationship of the
publishers or authors of the Document to the Document's overall subject
(or to related matters) and contains nothing that could fall directly
within that overall subject.  (Thus, if the Document is in part a
textbook of mathematics, a Secondary Section may not explain any
mathematics.)  The relationship could be a matter of historical
connection with the subject or with related matters, or of legal,
commercial, philosophical, ethical or political position regarding
them.

The \textbf{``Invariant Sections''} are certain Secondary Sections whose titles
are designated, as being those of Invariant Sections, in the notice
that says that the Document is released under this License.  If a
section does not fit the above definition of Secondary then it is not
allowed to be designated as Invariant.  The Document may contain zero
Invariant Sections.  If the Document does not identify any Invariant
Sections then there are none.

The \textbf{``Cover Texts''} are certain short passages of text that are listed,
as Front-Cover Texts or Back-Cover Texts, in the notice that says that
the Document is released under this License.  A Front-Cover Text may
be at most 5 words, and a Back-Cover Text may be at most 25 words.

A \textbf{``Transparent''} copy of the Document means a machine-readable copy,
represented in a format whose specification is available to the
general public, that is suitable for revising the document
straightforwardly with generic text editors or (for images composed of
pixels) generic paint programs or (for drawings) some widely available
drawing editor, and that is suitable for input to text formatters or
for automatic translation to a variety of formats suitable for input
to text formatters.  A copy made in an otherwise Transparent file
format whose markup, or absence of markup, has been arranged to thwart
or discourage subsequent modification by readers is not Transparent.
An image format is not Transparent if used for any substantial amount
of text.  A copy that is not ``Transparent'' is called \textbf{``Opaque''}.

Examples of suitable formats for Transparent copies include plain
ASCII without markup, Texinfo input format, LaTeX input format, SGML
or XML using a publicly available DTD, and standard-conforming simple
HTML, PostScript or PDF designed for human modification.  Examples of
transparent image formats include PNG, XCF and JPG.  Opaque formats
include proprietary formats that can be read and edited only by
proprietary word processors, SGML or XML for which the DTD and/or
processing tools are not generally available, and the
machine-generated HTML, PostScript or PDF produced by some word
processors for output purposes only.

The \textbf{``Title Page''} means, for a printed book, the title page itself,
plus such following pages as are needed to hold, legibly, the material
this License requires to appear in the title page.  For works in
formats which do not have any title page as such, ``Title Page'' means
the text near the most prominent appearance of the work's title,
preceding the beginning of the body of the text.

A section \textbf{``Entitled XYZ''} means a named subunit of the Document whose
title either is precisely XYZ or contains XYZ in parentheses following
text that translates XYZ in another language.  (Here XYZ stands for a
specific section name mentioned below, such as \textbf{``Acknowledgements''},
\textbf{``Dedications''}, \textbf{``Endorsements''}, or \textbf{``History''}.)  
To \textbf{``Preserve the Title''}
of such a section when you modify the Document means that it remains a
section ``Entitled XYZ'' according to this definition.

The Document may include Warranty Disclaimers next to the notice which
states that this License applies to the Document.  These Warranty
Disclaimers are considered to be included by reference in this
License, but only as regards disclaiming warranties: any other
implication that these Warranty Disclaimers may have is void and has
no effect on the meaning of this License.


\begin{center}
{\Large\bf 2. VERBATIM COPYING}
\end{center}

You may copy and distribute the Document in any medium, either
commercially or noncommercially, provided that this License, the
copyright notices, and the license notice saying this License applies
to the Document are reproduced in all copies, and that you add no other
conditions whatsoever to those of this License.  You may not use
technical measures to obstruct or control the reading or further
copying of the copies you make or distribute.  However, you may accept
compensation in exchange for copies.  If you distribute a large enough
number of copies you must also follow the conditions in section 3.

You may also lend copies, under the same conditions stated above, and
you may publicly display copies.


\begin{center}
{\Large\bf 3. COPYING IN QUANTITY}
\end{center}


If you publish printed copies (or copies in media that commonly have
printed covers) of the Document, numbering more than 100, and the
Document's license notice requires Cover Texts, you must enclose the
copies in covers that carry, clearly and legibly, all these Cover
Texts: Front-Cover Texts on the front cover, and Back-Cover Texts on
the back cover.  Both covers must also clearly and legibly identify
you as the publisher of these copies.  The front cover must present
the full title with all words of the title equally prominent and
visible.  You may add other material on the covers in addition.
Copying with changes limited to the covers, as long as they preserve
the title of the Document and satisfy these conditions, can be treated
as verbatim copying in other respects.

If the required texts for either cover are too voluminous to fit
legibly, you should put the first ones listed (as many as fit
reasonably) on the actual cover, and continue the rest onto adjacent
pages.

If you publish or distribute Opaque copies of the Document numbering
more than 100, you must either include a machine-readable Transparent
copy along with each Opaque copy, or state in or with each Opaque copy
a computer-network location from which the general network-using
public has access to download using public-standard network protocols
a complete Transparent copy of the Document, free of added material.
If you use the latter option, you must take reasonably prudent steps,
when you begin distribution of Opaque copies in quantity, to ensure
that this Transparent copy will remain thus accessible at the stated
location until at least one year after the last time you distribute an
Opaque copy (directly or through your agents or retailers) of that
edition to the public.

It is requested, but not required, that you contact the authors of the
Document well before redistributing any large number of copies, to give
them a chance to provide you with an updated version of the Document.


\begin{center}
{\Large\bf 4. MODIFICATIONS}
\end{center}

You may copy and distribute a Modified Version of the Document under
the conditions of sections 2 and 3 above, provided that you release
the Modified Version under precisely this License, with the Modified
Version filling the role of the Document, thus licensing distribution
and modification of the Modified Version to whoever possesses a copy
of it.  In addition, you must do these things in the Modified Version:

\begin{itemize}
\item[A.] 
   Use in the Title Page (and on the covers, if any) a title distinct
   from that of the Document, and from those of previous versions
   (which should, if there were any, be listed in the History section
   of the Document).  You may use the same title as a previous version
   if the original publisher of that version gives permission.
   
\item[B.]
   List on the Title Page, as authors, one or more persons or entities
   responsible for authorship of the modifications in the Modified
   Version, together with at least five of the principal authors of the
   Document (all of its principal authors, if it has fewer than five),
   unless they release you from this requirement.
   
\item[C.]
   State on the Title page the name of the publisher of the
   Modified Version, as the publisher.
   
\item[D.]
   Preserve all the copyright notices of the Document.
   
\item[E.]
   Add an appropriate copyright notice for your modifications
   adjacent to the other copyright notices.
   
\item[F.]
   Include, immediately after the copyright notices, a license notice
   giving the public permission to use the Modified Version under the
   terms of this License, in the form shown in the Addendum below.
   
\item[G.]
   Preserve in that license notice the full lists of Invariant Sections
   and required Cover Texts given in the Document's license notice.
   
\item[H.]
   Include an unaltered copy of this License.
   
\item[I.]
   Preserve the section Entitled ``History'', Preserve its Title, and add
   to it an item stating at least the title, year, new authors, and
   publisher of the Modified Version as given on the Title Page.  If
   there is no section Entitled ``History'' in the Document, create one
   stating the title, year, authors, and publisher of the Document as
   given on its Title Page, then add an item describing the Modified
   Version as stated in the previous sentence.
   
\item[J.]
   Preserve the network location, if any, given in the Document for
   public access to a Transparent copy of the Document, and likewise
   the network locations given in the Document for previous versions
   it was based on.  These may be placed in the ``History'' section.
   You may omit a network location for a work that was published at
   least four years before the Document itself, or if the original
   publisher of the version it refers to gives permission.
   
\item[K.]
   For any section Entitled ``Acknowledgements'' or ``Dedications'',
   Preserve the Title of the section, and preserve in the section all
   the substance and tone of each of the contributor acknowledgements
   and/or dedications given therein.
   
\item[L.]
   Preserve all the Invariant Sections of the Document,
   unaltered in their text and in their titles.  Section numbers
   or the equivalent are not considered part of the section titles.
   
\item[M.]
   Delete any section Entitled ``Endorsements''.  Such a section
   may not be included in the Modified Version.
   
\item[N.]
   Do not retitle any existing section to be Entitled ``Endorsements''
   or to conflict in title with any Invariant Section.
   
\item[O.]
   Preserve any Warranty Disclaimers.
\end{itemize}

If the Modified Version includes new front-matter sections or
appendices that qualify as Secondary Sections and contain no material
copied from the Document, you may at your option designate some or all
of these sections as invariant.  To do this, add their titles to the
list of Invariant Sections in the Modified Version's license notice.
These titles must be distinct from any other section titles.

You may add a section Entitled ``Endorsements'', provided it contains
nothing but endorsements of your Modified Version by various
parties--for example, statements of peer review or that the text has
been approved by an organization as the authoritative definition of a
standard.

You may add a passage of up to five words as a Front-Cover Text, and a
passage of up to 25 words as a Back-Cover Text, to the end of the list
of Cover Texts in the Modified Version.  Only one passage of
Front-Cover Text and one of Back-Cover Text may be added by (or
through arrangements made by) any one entity.  If the Document already
includes a cover text for the same cover, previously added by you or
by arrangement made by the same entity you are acting on behalf of,
you may not add another; but you may replace the old one, on explicit
permission from the previous publisher that added the old one.

The author(s) and publisher(s) of the Document do not by this License
give permission to use their names for publicity for or to assert or
imply endorsement of any Modified Version.


\begin{center}
{\Large\bf 5. COMBINING DOCUMENTS}
\end{center}


You may combine the Document with other documents released under this
License, under the terms defined in section 4 above for modified
versions, provided that you include in the combination all of the
Invariant Sections of all of the original documents, unmodified, and
list them all as Invariant Sections of your combined work in its
license notice, and that you preserve all their Warranty Disclaimers.

The combined work need only contain one copy of this License, and
multiple identical Invariant Sections may be replaced with a single
copy.  If there are multiple Invariant Sections with the same name but
different contents, make the title of each such section unique by
adding at the end of it, in parentheses, the name of the original
author or publisher of that section if known, or else a unique number.
Make the same adjustment to the section titles in the list of
Invariant Sections in the license notice of the combined work.

In the combination, you must combine any sections Entitled ``History''
in the various original documents, forming one section Entitled
``History''; likewise combine any sections Entitled ``Acknowledgements'',
and any sections Entitled ``Dedications''.  You must delete all sections
Entitled ``Endorsements''.

\begin{center}
{\Large\bf 6. COLLECTIONS OF DOCUMENTS}
\end{center}

You may make a collection consisting of the Document and other documents
released under this License, and replace the individual copies of this
License in the various documents with a single copy that is included in
the collection, provided that you follow the rules of this License for
verbatim copying of each of the documents in all other respects.

You may extract a single document from such a collection, and distribute
it individually under this License, provided you insert a copy of this
License into the extracted document, and follow this License in all
other respects regarding verbatim copying of that document.


\begin{center}
{\Large\bf 7. AGGREGATION WITH INDEPENDENT WORKS}
\end{center}


A compilation of the Document or its derivatives with other separate
and independent documents or works, in or on a volume of a storage or
distribution medium, is called an ``aggregate'' if the copyright
resulting from the compilation is not used to limit the legal rights
of the compilation's users beyond what the individual works permit.
When the Document is included in an aggregate, this License does not
apply to the other works in the aggregate which are not themselves
derivative works of the Document.

If the Cover Text requirement of section 3 is applicable to these
copies of the Document, then if the Document is less than one half of
the entire aggregate, the Document's Cover Texts may be placed on
covers that bracket the Document within the aggregate, or the
electronic equivalent of covers if the Document is in electronic form.
Otherwise they must appear on printed covers that bracket the whole
aggregate.


\begin{center}
{\Large\bf 8. TRANSLATION}
\end{center}


Translation is considered a kind of modification, so you may
distribute translations of the Document under the terms of section 4.
Replacing Invariant Sections with translations requires special
permission from their copyright holders, but you may include
translations of some or all Invariant Sections in addition to the
original versions of these Invariant Sections.  You may include a
translation of this License, and all the license notices in the
Document, and any Warranty Disclaimers, provided that you also include
the original English version of this License and the original versions
of those notices and disclaimers.  In case of a disagreement between
the translation and the original version of this License or a notice
or disclaimer, the original version will prevail.

If a section in the Document is Entitled ``Acknowledgements'',
``Dedications'', or ``History'', the requirement (section 4) to Preserve
its Title (section 1) will typically require changing the actual
title.


\begin{center}
{\Large\bf 9. TERMINATION}
\end{center}


You may not copy, modify, sublicense, or distribute the Document except
as expressly provided for under this License.  Any other attempt to
copy, modify, sublicense or distribute the Document is void, and will
automatically terminate your rights under this License.  However,
parties who have received copies, or rights, from you under this
License will not have their licenses terminated so long as such
parties remain in full compliance.


\begin{center}
{\Large\bf 10. FUTURE REVISIONS OF THIS LICENSE}
\end{center}


The Free Software Foundation may publish new, revised versions
of the GNU Free Documentation License from time to time.  Such new
versions will be similar in spirit to the present version, but may
differ in detail to address new problems or concerns.  See
http://www.gnu.org/copyleft/.

Each version of the License is given a distinguishing version number.
If the Document specifies that a particular numbered version of this
License ``or any later version'' applies to it, you have the option of
following the terms and conditions either of that specified version or
of any later version that has been published (not as a draft) by the
Free Software Foundation.  If the Document does not specify a version
number of this License, you may choose any version ever published (not
as a draft) by the Free Software Foundation.

%---------------------------------------------------------------------

%%% Local Variables: 
%%% mode: latex
%%% TeX-master: "main"
%%% End: 





\bibliographystyle{is-alpha}
\bibliography{bibliografia,rfc}



\end{document}

%%% Local Variables: 
%%% mode: latex
%%% TeX-master: t
%%% End: 
